\documentclass[11pt]{article}
\usepackage{fontspec}
\usepackage{xeCJK}
\usepackage{amsmath}
\usepackage{amssymb}
\usepackage{array}
\usepackage{booktabs}
\usepackage[margin=1in]{geometry}
\usepackage[hidelinks]{hyperref}
\usepackage{seqsplit}

\setmainfont{Times New Roman}
\setCJKmainfont{Malgun Gothic}
\setCJKsansfont{Malgun Gothic}
\setCJKmonofont{Malgun Gothic}
\xeCJKsetup{CJKspace=true}
\XeTeXlinebreaklocale "ko"
\XeTeXlinebreakskip=0pt plus 1pt
\setlength{\emergencystretch}{2em}

\newcommand{\claimstatus}[2]{\textbf{#1 [#2]}}
\newcommand{\poly}{\operatorname{poly}}
\newcommand{\bits}{\operatorname{bitlength}}
\newcommand{\archivepaper}{\path{paper/main-ko.tex}}
\newcolumntype{L}[1]{>{\raggedright\arraybackslash}p{#1}}
\newenvironment{claimproof}{\par\noindent\textbf{증명.}}{\hfill\(\square\)\par}
\newenvironment{costmodel}[1]{%
  \begin{center}
  \small
  \begin{tabular}{L{0.19\linewidth}L{0.69\linewidth}}
  \toprule
  \multicolumn{2}{l}{\textbf{#1}}\\
  \midrule
}{%
  \bottomrule
  \end{tabular}
  \end{center}
}
\newcommand{\costrow}[2]{\textbf{#1} & #2\tabularnewline}
\renewcommand{\abstractname}{초록}


\title{인식되지 않는 유전적 차수 약속 아래의\\
Las Vegas 완전 인수분해}
\author{MOSEF 연구 프로그램}
\date{2026년 7월 31일}

\begin{document}
\maketitle

\begin{abstract}
고전적 exponent--GCD 인수분해의 성공 조건은 schedule이 차수를 단순히
포함하는지가 아니라 서로 다른 소인수 성분에 서로 다른 capped
divisibility profile을 부여하는지이다. 이 separation 의미론을 바탕으로
본 논문은 hereditary order-asymmetry 약속 아래의 두 완전 인수분해
정리를 제시한다. 첫 번째 알고리즘은 \(p-1\) schedule의 각 시도에서
새로운 균등 밑을 뽑고 완전한 witness cycle마다 적어도 \(5/12\)의
확률로 분할한다. 두 번째 알고리즘은 비분할 \(p+1\) schedule에서
새로운 균등 Lucas 매개변수를 사용하며 witness 시도마다 \(1/12\)보다
큰 성공 확률을 갖는다. 정확한 소수 판정, 최대 perfect-power 전처리,
모든 modular 평가와 GCD, 검증된 재귀 및 출력 비용은 입력의 표준 이진
길이에 과금한다. witness 소인수와 promise membership은 online
알고리즘에 주어지지 않으며 알고리즘도 이를 인식하지 않는다. 노드별
유한 예산 \(s\)를 둔 total wrapper는 모든 입력에서 완전한 검증
인수분해 또는 \texttt{UNRESOLVED}를 반환한다. 두 약속 안의 지역
미해결 꼬리는 각각 \((7/12)^s\), \((11/12)^s\) 이하이고, 재귀 노드의
독립성을 가정하지 않은 전체 재귀 상한은 각각
\(\min\{1,4m(7/12)^s\}\), \(\min\{1,4m(11/12)^s\}\)이다. 고정 밑에는
명시적 hereditary 반례가 있으므로 fresh sampling이 필수이다. 또한
divisor coverage와 order-signature injectivity를 구별하고, 켤레 Lucas
매개변수가 독립 채널을 만들지 않음을 보이며, 공통 schedule의 유한
분포 한계를 기록한다. 증명되지 않은 수론적 추측은 가정하지 않지만
이는 인식되지 않는 약속 입력의 결과이지 임의 정수에 대한 고전적
다항시간 인수분해 알고리즘이 아니다. 정리의 성공 확률과 total
wrapper의 미해결 상한은 서로 다른 보장으로 분리한다.
\end{abstract}

\begin{costmodel}{비용 모델 요약}
\costrow{온라인 과금}{공개 schedule 구성, 정확한 전처리, fresh sampling,
modular \(p-1\) 또는 Lucas 평가, 모든 GCD와 약수 검증, 재귀, 검증된
출력을 포함한다.}
\costrow{주어지지 않음}{witness 소인수, promise-membership bit,
factor-dependent schedule은 입력이 아니며 추론하지도 않는다.}
\costrow{오프라인 증명}{군 및 root 개수, tail 상한, 반례 탐색,
재현 기록은 정리를 뒷받침하지만 online advice가 아니다.}
\costrow{보장 범위}{기대 다항 bit 복잡도는 명시한 hereditary 약속에서
성립한다. bounded wrapper는 모든 입력에서 종료하지만
\texttt{UNRESOLVED}를 반환할 수 있다.}
\end{costmodel}

\section{범위와 계산 모형}

\(N\ge1\)에 대해
\[
m=\bits(N)=\lfloor\log_2N\rfloor+1
\]
로 둔다. 입력은 이진수이며 모든 복잡도 진술은 각 재귀 노드에서 앞의
비용 표를 적용한다. 정확한 소수 판정과 최대 완전거듭제곱 분해를
splitter보다 먼저 수행하고, 보고된 약수는 두 자식으로 재귀하기 전에
검증한다.

\(B(k),R(k)\)를 다항식으로 유계이고 다항시간에 계산되는 고정된
비감소 양의 정수 함수라 하고
\[
M_B(k)=\operatorname{lcm}(1,\ldots,B(k))
\]
로 둔다. 지역 비트 길이가 \(k\)인 합성 비완전거듭제곱 \(K\)에 대해,
서로 다른 \(p,q\mid K\)와 \(1\le t\le R(k)\)가
\[
p-1\mid tM_B(k),\qquad q-1\nmid tM_B(k)
\]
를 만족하면 \(p-1\) witness라 한다. 비분할 \(p+1\) witness는 홀수
잔여 노드에서 \(p+1\mid tM_B(k)\), \(q+1\nmid tM_B(k)\)로 정의한다.
재귀 중 도달하는 모든 합성 비완전거듭제곱 잔여 노드에 witness가
존재할 때 약속은 유전적이다. 이 정의는 미지의 소인수에 의존하므로
온라인 알고리즘은 witness를 받지 않으며 membership을 판정하지도 않는다.

\section{새로운 밑을 이용한 \texorpdfstring{\(p-1\)}{p-1} 인수분해}

\claimstatus{THM-001}{PROVED}. 위의 고정된 \(B,R\)마다, 새로운 균등
잉여류를 사용하는 고전적 Las Vegas 알고리즘은 유전적 \(p-1\)
semismooth-asymmetry 약속의 모든 입력을 완전히 인수분해한다. 알고리즘은
항상 옳고 거의 확실하게 종료하며 기대 비트 복잡도는 입력 길이의
다항식이다. 약속된 잔여 노드의 한 완전 schedule cycle은 적어도
\(5/12\)의 확률로 분할한다.

\begin{claimproof}
잔여 노드 \(K\)에서 \(1\le t\le R(k)\)를 순회한다. 각 \(t\)마다 새로운
정확한 균등 \(a\bmod K\)를 뽑아 먼저 \(\gcd(a,K)\)를 검사하고, 이어서
\[
\gcd\!\left(a^{tM_B(k)}-1,K\right)
\]
를 계산한다. witness \(p,q,t\)를 고정하면 모든 unit은 \(p\)에서 1이
되지만, \(q\)에서 지수에 의해 소거되는 unit은 진부분군이므로 전체
unit의 절반 이하이다. proper nonunit은 이미 약수를 준다. 따라서
\[
\Pr(\text{성공})\ge
\frac{K-1-\varphi(K)/2}{K}
\ge\frac{K-1}{2K}\ge\frac5{12}.
\]
새로운 표본은 기하 꼬리와 거의 확실한 종료를 준다. 지수 비트 길이는
\(O(B(k)\log B(k)+\log R(k))\)이고 인수분해 재귀 트리의 노드 수는
\(O(m)\)이므로 기대 비용은 \(\poly(m)\)이다. 완전한 branch 및 재귀
증명은 \path{research/proofs/THM-001-semismooth-promise.md}에 있다.
\end{claimproof}

고정 밑으로 바꾸면 정리가 깨진다. \(N=51\), \(a=2\),
\(d=\operatorname{lcm}(1,\ldots,8)=840\)에서 \(3-1\mid d\),
\(17-1\nmid d\)이지만 \(\operatorname{ord}_{17}(2)=8\mid d\)이므로
\(\gcd(2^{840}-1,51)=51\)인 full collision이 생긴다.

\section{포함은 분리가 아니다}

유한 양의 지수족 \(\Delta\)에 대해
\[
\Sigma_\Delta(r)=\{d\in\Delta:r\mid d\}
\]
를 정의한다.

\claimstatus{BAR-001}{PROVED}. 약수 포함은 차수 분리를 함의하지 않는다.
profile \((r_1,\ldots,r_s)\)는 그 signature들이 모두 같지 않을 때와
그때에만 분리된다. \(\{1,\ldots,n\}\)의 서로 다른 모든 차수 쌍을
분리하는 필요충분조건은 \(r\mapsto\Sigma_\Delta(r)\)의 단사성이며,
비어 있지 않은 명시적 signature족은
\[
|\Delta|\ge\lceil\log_2(n+1)\rceil
\]
을 만족한다.

\begin{claimproof}
한 지수가 일부 \(r_i\)만 나누는 조건은 그 지수가 일부 signature에만
속하는 조건과 같다. 따라서 전체 분리는 signature 불일치와 동치이다.
비어 있지 않은 부분집합은 최대 \(2^{|\Delta|}-1\)개이므로 하한이
따른다. 최소 반례 \(\Delta=\{2\}\)에서는 차수 1과 2가 모두 포함되지만
signature가 같고 \(\gcd(5^2-1,6)=6\)이다.
\end{claimproof}

\section{켤레 채널의 상관}

\(V_0(P,1)=2\), \(V_1(P,1)=P\),
\(V_{j+1}(P,1)=PV_j(P,1)-V_{j-1}(P,1)\)로 둔다.

\claimstatus{BAR-002}{PROVED}. \(a\)가 \(N\)의 unit이고
\(P=a+a^{-1}\)이면 모든 \(d\ge0\)에 대해
\[
V_d(P,1)-2=a^{-d}(a^d-1)^2,\qquad
P^2-4=a^{-2}(a^2-1)^2\pmod N.
\]
따라서 유도된 Lucas 잉여와 곱셈 잉여는 같은 소수 support를 갖는다.
지수 2를 보존하면 이 유도 채널은 곱셈족의 성공 영역을 넓히지 못한다.

\begin{claimproof}
\(a^d+a^{-d}\)가 같은 초기값과 점화식을 만족하므로
\(V_d(P,1)\)와 같다. 2를 빼고 unit \(a^d\)를 곱하면 제곱 항등식이
나온다. 판별식 항등식은 \(d=2\)인 경우이다. 소수거듭제곱에서는
valuation이 두 배가 되어 오히려 분할을 악화시킬 수 있다.
\((N,a,P,d)=(25,2,15,4)\)에서 곱셈 GCD는 5이지만 Lucas GCD는 25이다.
\end{claimproof}

\section{독립적인 비분할 \texorpdfstring{\(p+1\)}{p+1} 인수분해}

\claimstatus{LEM-003}{PROVED}. 모든 홀수 소수 \(q\)와 \(d>0\)에 대해
\[
\#\{P\in\mathbb F_q:V_d(P,1)=2\}
=\frac{\gcd(d,q-1)+\gcd(d,q+1)}2.
\]
이 중 정확히 \(1+\mathbf1_{2\mid d}\)개가 퇴화 매개변수 \(P=\pm2\)이고,
정확히 \((q-1)/2\)개가 비제곱 판별식을 갖는다.

\begin{claimproof}
\(P=z+z^{-1}\)로 매개화하고 \(z\)를 차수 \(q-1\)인
\(\mathbb F_q^\times\)와 차수 \(q+1\)인
\(\mathbb F_{q^2}^\times\)의 norm-one 부분군으로 나눈다.
\(V_d(P,1)=2\)는 \(z^d=1\)과 동치이다. 비퇴화 쌍
\(\{z,z^{-1}\}\)을 한 번씩 세고 \(z=\pm1\) 고정점을 보정하면 식이
따른다.
\end{claimproof}

\claimstatus{THM-002}{PROVED}. 새로운 균등 \(P\bmod K\)를 사용하는
고전적 Las Vegas 알고리즘은 유전적 비분할 \(p+1\)-asymmetry 약속의
모든 입력을 완전히 인수분해한다. 항상 옳고 거의 확실하게 종료하며
기대 비트 복잡도는 다항식이다. witness 시도의 성공 확률은
\(1/12\)보다 크다.

\begin{claimproof}
\(\gcd(P^2-4,K)\)의 proper factor를 보존하고, 이 GCD가 full이어도
\(\gcd(V_d(P,1)-2,K)\)를 별도로 검사한다. witness 소수 \(p\)에서는
\((p-1)/2\)개의 비분할 매개변수가 수열 값을 2로 만든다. \(q\)에서
비퇴화 collision 수는 LEM-003에 의해
\[
C_q(d)\le(3q-9)/4
\]
이다. CRT 사건의 확률은
\[
\frac{p-1}{2p}\frac{q-C_q(d)}q>\frac1{12}
\]
이다. 새로운 매개변수, 이진 Lucas 평가, 전처리와 재귀는 기하 꼬리와
다항 기대 비용을 준다. 완전한 증명은
\path{research/proofs/THM-002-nonsplit-lucas-promise.md}에 있다.
\end{claimproof}

\section{유한 예산 total wrapper}

양의 정수 \(s\)를 고정하고 각 randomized composite 노드의 무한 cycle
반복을 최대 \(s\)개의 완전 cycle로 바꾼다. 소수, 최대
완전거듭제곱, 짝수 인수, proper factor, 예산 소진 분기를 모두
명시한다. 공개 결과는 입력의 완전한 소인수 목록 또는
\texttt{UNRESOLVED}뿐이며 부분 인수 목록은 폐기한다.

이 bounded 형태는 THM-001과 THM-002의 total corollary이다. 모든 양의
입력에서 종료하고 잘못된 인수분해를 반환하지 않는다. 각 유전적 약속
안에서 randomized 노드에 도달했다는 조건 아래 지역 미해결 확률은
\[
\Pr[U_{p-1}]\le\left(\frac7{12}\right)^s,\qquad
\Pr[U_{p+1}]\le\left(\frac{11}{12}\right)^s
\]
이다. 최대거듭제곱 전처리 뒤 전체 재귀 호출 수가 \(4m\)보다 작으므로
재귀 노드 사이의 독립성을 가정하지 않는 조건부 union bound는
\[
\Pr[U^{\mathrm{complete}}_{p-1}]
\le\min\left\{1,4m\left(\frac7{12}\right)^s\right\},
\]
\[
\Pr[U^{\mathrm{complete}}_{p+1}]
\le\min\left\{1,4m\left(\frac{11}{12}\right)^s\right\}
\]
를 준다. 약속 밖에서는 totality와 no-wrong-factor 정확성만 남는다.
\texttt{UNRESOLVED}는 membership 판정이 아니다. 전체 증명과 실행
의미론은 다음 파일에 있다.
\par\smallskip
\noindent\begin{minipage}{\linewidth}
\raggedright\small
\path{research/proofs/M84-bounded-total-promise-wrappers.md}\\
\path{python/mosef_reference/promise_wrappers.py}
\end{minipage}

\section{공통 schedule의 한계}

\[
\sigma_\Delta(p)=
\bigl(\mathbf1_{p-1\mid d},\mathbf1_{p+1\mid d}\bigr)_{d\in\Delta}
\]
는 미지 소인수에서 평가되는 분석 객체이며 입력 recognizer가 아니다.

\claimstatus{BAR-003}{PROVED}. 서로 다른 홀수 소수 \(p,q\)에 대해 지역
\(p-1\), \(p+1\) 비대칭 약속의 합집합은
\(\sigma_\Delta(p)\ne\sigma_\Delta(q)\)와 동치이다. 유한 소수집합
\(S\)에서 \(s=|S|\), nonzero signature 소수 수를 \(h\),
\(D=\max\Delta\)라 하면 약속 쌍 비율은
\[
\rho_\Delta(S)\le
\frac{h(2s-h-1)}{s(s-1)}
\le\min\left\{1,\frac{8|\Delta|\sqrt D}{s-1}\right\}.
\]

\begin{claimproof}
서로 다른 signature는 정확히 한 소수 쪽만 포함하는 좌표를 준다.
nonzero signature에 닿는 쌍을 세면 첫 부등식이 나온다. 각 \(d\)에
대해 \(p-1\mid d\) 또는 \(p+1\mid d\)인 소수 수는 약수 수로 유계되고
\(\tau(d)\le2\sqrt d\)이므로 \(h\le4|\Delta|\sqrt D\)이다.
\end{claimproof}

\claimstatus{BAR-004}{PROVED}. 최대 지수 비트 길이가
\(o(k\log k)\)인 다항 크기 명시적 지수 목록은 \(p\pm1\) 약수 검사를
통해 최대 \(2^{o(k)}\)개의 홀수 소수를 hit한다. 별도로 가정된
\(2^{\alpha k}\) 크기의 공통 길이 소수 모집단에서는 covered-pair
비율이 0으로 수렴한다.

\begin{claimproof}
모든 hit은 \(p-1\mid d\) 또는 \(p+1\mid d\)에서 생기므로 전체 수는
지수들의 약수 수 합의 두 배 이하이다.
\(\log_2\tau(d)=O(\ell(d)/\log\ell(d))=o(k)\)를 다항 개의 지수에
합쳐도 \(2^{o(k)}\)이다. BAR-003의 정확한 쌍 경계에 대입하면 된다.
이는 모든 compact 표현이나 adaptive schedule의 하한이 아니다.
\end{claimproof}

\section{1차 문헌 포지셔닝}

\paragraph{M83-R01: Pollard \(p-1\).}
Pollard의 출판 절차는 unit 및 full-collision 분기를 포함한 고전적
smooth-\(p-1\) 메커니즘을 제공한다 \cite{pollard1974theorems}. 그러나
유전적 약속, 매 시도의 fresh exact-uniform sampling, \(5/12\) 완전
cycle 하한, THM-001의 완결된 완전 재귀 비용은 진술하지 않는다.
따라서 THM-001은 확립된 메커니즘 주위의 제한적 종합으로 제시하며
우선권을 주장하지 않는다.

\paragraph{M83-R02: Williams \(p+1\).}
Williams는 Lucas sequence 메커니즘과 split, nonsplit,
zero-discriminant 가정을 제공한다 \cite{williams1982pplusone}.
LEM-003의 정확한 매개변수 root count, 모든 퇴화 분기의 처리,
THM-002의 \(1/12\)보다 큰 witness 하한은 저장소의 자기완결
유도이다. 이 감사는 이에 문헌 우선권 지위를 부여하지 않는다.

\paragraph{M83-R03: 분리 signature.}
이진 membership signature로 모든 쌍을 구분하는 언어는 고전적
separating-system 개념이다 \cite{katona1966separating}. BAR-001과
BAR-024는 이를 divisor/support 좌표, 소수거듭제곱 valuation, 정확한
proper-GCD 의미론에 특수화한다. 이 특수화가 논문의 범위이며, 이진
signature 단사성 자체를 새로운 조합론 개념으로 주장하지 않는다.

등록된 M83 행렬은 전체 1차 문헌 비교를 보존한다. 이 제한 감사에서
동일한 정식화를 찾지 못한 사실은 독창성이나 우선권의 증거가 아니다.

\section{재현성과 한계}

\begin{verbatim}
python scripts/run_m3_semismooth_search.py
python scripts/check_m3_semismooth_differential.py
python scripts/run_m7_nonsplit_search.py
python scripts/check_m7_nonsplit_differential.py
python scripts/check_m84_promise_wrappers.py
python -m unittest discover -s tests -p test_promise_wrappers.py
\end{verbatim}
EXP-0003과 EXP-0006의 실험 기록은 \path{research/experiments}에 있다.
실험은 유한 산술과 branch를 검사하며 확률 하한의 증명을 대체하지 않는다.

archival monograph는 전체 증명, 반례, 비용 장부, 부정 결과, claim
의존성을 보존하고 M82 projection은 집중 논문 배치를 검증한다. 약속
membership은 여전히 인식되지 않고 입력 밀도 정리도 없으며, 일반
고전적 다항시간 인수분해 문제는 열린 상태이다.

\bibliographystyle{plain}
\small
\bibliography{paper/references}

\end{document}
